%% bare_conf.tex
%% V1.4
%% 2012/12/27
%% by Michael Shell
%% See:
%% http://www.michaelshell.org/
%% for current contact information.
%%
%% This is a skeleton file demonstrating the use of IEEEtran.cls
%% (requires IEEEtran.cls version 1.8 or later) with an IEEE conference paper.
%%
%% Support sites:
%% http://www.michaelshell.org/tex/ieeetran/
%% http://www.ctan.org/tex-archive/macros/latex/contrib/IEEEtran/
%% and
%% http://www.ieee.org/

%%*************************************************************************
%% Legal Notice:
%% This code is offered as-is without any warranty either expressed or
%% implied; without even the implied warranty of MERCHANTABILITY or
%% FITNESS FOR A PARTICULAR PURPOSE! 
%% User assumes all risk.
%% In no event shall IEEE or any contributor to this code be liable for
%% any damages or losses, including, but not limited to, incidental,
%% consequential, or any other damages, resulting from the use or misuse
%% of any information contained here.
%%
%% All comments are the opinions of their respective authors and are not
%% necessarily endorsed by the IEEE.
%%
%% This work is distributed under the LaTeX Project Public License (LPPL)
%% ( http://www.latex-project.org/ ) version 1.3, and may be freely used,
%% distributed and modified. A copy of the LPPL, version 1.3, is included
%% in the base LaTeX documentation of all distributions of LaTeX released
%% 2003/12/01 or later.
%% Retain all contribution notices and credits.
%% ** Modified files should be clearly indicated as such, including  **
%% ** renaming them and changing author support contact information. **
%%
%% File list of work: IEEEtran.cls, IEEEtran_HOWTO.pdf, bare_adv.tex,
%%                    bare_conf.tex, bare_jrnl.tex, bare_jrnl_compsoc.tex,
%%                    bare_jrnl_transmag.tex
%%*************************************************************************

\documentclass[conference]{IEEEtran}

\usepackage[utf8]{inputenc}
\usepackage[T1]{fontenc}
\usepackage[english]{babel}
\usepackage{amsmath,amssymb}
\usepackage{bm}
\usepackage{booktabs}
\usepackage{array}
\usepackage{xcolor}
\usepackage{hyperref}
\usepackage{float}
\usepackage{enumitem}
\usepackage{listings}
\usepackage{caption}
\usepackage{subcaption}
\usepackage[numbers,sort&compress]{natbib}
\usepackage{rotating}
\usepackage{graphicx}
\graphicspath{{../notebooks/figures/}}

\hypersetup{colorlinks=true,linkcolor=blue!70!black,
  citecolor=green!50!black,urlcolor=blue!70!black}

\definecolor{codebg}{RGB}{248,248,248}
\lstdefinestyle{pythonstyle}{
  backgroundcolor=\color{codebg},basicstyle=\ttfamily\small,
  keywordstyle=\bfseries,frame=single,breaklines=true,
  language=Python,showstringspaces=false,
  numbers=left,numberstyle=\tiny\color{gray},xleftmargin=2em}
\lstset{style=pythonstyle}

\newcommand{\btheta}{\bm{\theta}}
\newcommand{\cD}{\mathcal{D}}
\newcommand{\cN}{\mathcal{N}}
\newcommand{\KL}{\mathrm{KL}}
\newcommand{\ELBO}{\mathcal{L}_{\text{ELBO}}}
\newcommand{\MAP}{\text{MAP}}
\newcolumntype{C}[1]{>{\centering\arraybackslash}p{#1}}
\newcolumntype{L}[1]{>{\raggedright\arraybackslash}p{#1}}
\newcommand{\todo}[1]{\noindent\colorbox{yellow!40}{\parbox{\dimexpr\linewidth-2\fboxsep}{[TODO] #1}}}

\begin{document}

\title{Bayesian Neural Networks:\\
\texttt{Implementations using Python Libraries}}

\author{\IEEEauthorblockN{Anonymous Author}
\IEEEauthorblockA{Department of Computer Science\\
University of Example\\
City, Country\\
Email: author@example.com}}

\maketitle

\section{Introduction}
\label{sec:intro}

Standard deep neural networks are notoriously miscalibrated: high confidence
predictions are not guaranteed to be correct, and models frequently fail
silently under distribution shift~\citep{goan2020bayesian,murphy2023pml}.
Bayesian Deep Learning addresses this by replacing a deterministic weight
estimate with a posterior distribution $p(\btheta \mid \cD)$, propagating weight
uncertainty into calibrated predictive distributions.
This report compares two PyTorch libraries implementing complementary paradigms:
\begin{itemize}[leftmargin=*]
  \item \texttt{bayesian-torch}~\citep{intel2021bayesiantorch}, a
        Bayesian by design library from Intel Labs implementing mean field
        Variational Inference via Flipout~\citep{wen2018flipout} and
        Bayes by Backprop~\citep{blundell2015weight}.
  \item \texttt{laplace-torch}~\citep{daxberger2021laplace,aleximmer2024laplace},
        a post hoc library applying the Laplace Approximation to any
        pre trained model without modifying its training loop.
\end{itemize}

Experiments span three settings: a synthetic binary classification
benchmark, a toy regression
experiment, and a replication of the
Laplace Redux PovertyMap experiment.
Full metric tables and figure galleries for completed experiments are in the
Appendix.

\section{Theoretical Background}
\label{sec:theory}

\subsection{Bayesian Deep Learning}

Given $\cD = \{(\bm{x}_n, y_n)\}_{n=1}^{N}$ and network $f_{\btheta}$, the
posterior is $p(\btheta \mid \cD) = p(\cD|\btheta)\,p(\btheta)/p(\cD)$.
The evidence is intractable for modern DNNs. With a Gaussian prior
$p(\btheta)=\cN(\bm{0},\tau^{-1}\mathbf{I})$, the MAP estimate
$\btheta_{\MAP} = \arg\max_{\btheta}[\log p(\cD|\btheta)+\log p(\btheta)]$
corresponds to $L_2$ regularised training with weight decay
$\lambda=\tau$~\citep{murphy2023pml}.

\subsection{Stochastic Variational Inference}

VI approximates the posterior with a tractable family $q_{\bm\phi}(\btheta)$
by maximising the Evidence Lower Bound~\citep{blundell2015weight,goan2020bayesian}:
\begin{equation}
  \ELBO(\bm\phi)
    = \mathbb{E}_{q}\!\bigl[\log p(\cD|\btheta)\bigr]
    - \KL\bigl(q_{\bm\phi}\,\|\,p(\btheta)\bigr).
\label{eq:elbo}
\end{equation}
A mean field Gaussian $q_{\bm\phi}(w_i) = \cN(\mu_i,\sigma_i^2)$ with
$\sigma_i = \log(1+\exp(\rho_i))>0$ keeps the KL in~\eqref{eq:elbo} in
closed form, but doubles the parameter count.

\subsubsection{Bayes by Backprop}

\citet{blundell2015weight} reparameterise each weight as
\begin{equation}
  w_i = \mu_i + \sigma_i\,\varepsilon_i, \qquad \varepsilon_i\sim\cN(0,1),
\end{equation}
enabling low variance gradient estimates of~\eqref{eq:elbo} via
backpropagation.
A single noise draw is shared across all examples in a
mini batch, which inflates gradient variance at large batch sizes.

\subsubsection{Flipout}

\citet{wen2018flipout} reduce this variance by applying independent sign flip
matrices $\mathbf{r}_n,\mathbf{s}_n\in\{-1,+1\}^*$ to the shared
perturbation $\Delta\hat{\mathbf{W}}$ for each example $n$:
\begin{equation}
  \hat{\mathbf{W}}_n
    = \mathbf{W}_{\MAP} + \mathbf{s}_n \odot \Delta\hat{\mathbf{W}}
      \odot \mathbf{r}_n^{\!\top}.
\end{equation}

\subsubsection{MOPED Initialisation}

Training mean field BNNs from random initialisations is fragile.
\texttt{bayesian-torch}~\citep{intel2021bayesiantorch} addresses this via
MOPED: the posterior means are warm started from a
pre trained MAP checkpoint and the prior scale is set proportionally:
\begin{equation}
  \mu_i^{(0)} = \theta_{\MAP}^{(i)}, \qquad
  \sigma_i^{\text{prior}} = \delta\,|\theta_{\MAP}^{(i)}|, \quad \delta=0.5.
\end{equation}

\subsubsection{AvUC Loss}
\label{sec:theory:avuc}

Standard ELBO training minimises NLL but does not explicitly calibrate the
relationship between confidence and accuracy.
The
Accuracy vs Uncertainty Calibration loss, implemented in
\texttt{bayesian-torch}~\citep{intel2021bayesiantorch}, adds a direct
penalty for miscalibrated confidence.

\paragraph{AvUC objective.}
The AvUC loss aggregates the fractions of samples in the two
bad quadrants~\citep{intel2021bayesiantorch}:
\begin{equation}
  \mathcal{L}_{\mathrm{AvUC}}
    = \underbrace{\frac{n_{AU}}{n_{AU}+n_{AC}}}_{\text{over uncertain among correct}}
    + \underbrace{\frac{n_{IC}}{n_{IC}+n_{IU}}}_{\text{overconfident among wrong}},
  \label{eq:avuc}
\end{equation}
where $n_X$ counts samples in quadrant $X$ within the current mini batch.
Each term lies in $[0,1]$; both are zero when calibration is perfect.

\subsection{Epistemic and Aleatoric Uncertainty Decomposition}
\label{sec:theory:unc_decomp}

For MC based BNNs, the total predictive uncertainty can be decomposed into
two orthogonal components~\citep{goan2020bayesian}:
\begin{align}
  \text{Total}   &= H\!\left[\frac{1}{S}\sum_{s=1}^{S} p(y|f_{\btheta_s}(\bm{x}))\right], \\
  \text{Aleatoric}&= \frac{1}{S}\sum_{s=1}^{S} H\!\left[p(y|f_{\btheta_s}(\bm{x}))\right], \\
  \text{Epistemic}&= \text{Total} - \text{Aleatoric}
    = I\!\left(y;\btheta\,|\,\bm{x},\cD\right),
  \label{eq:mi}
\end{align}

\subsection{Laplace Approximation}
\label{sec:theory:la}

The LA~\citep{daxberger2021laplace} builds a Gaussian posterior around
$\btheta_{\MAP}$ via second order Taylor expansion:
\begin{equation}
  p(\btheta|\cD) \approx \cN\!\bigl(\btheta;\,\btheta_{\MAP},\,\mathbf{H}^{-1}\bigr),
  \qquad
  \mathbf{H} = -\nabla^2\log p(\btheta_{\MAP}|\cD) \succ 0.
\end{equation}

\paragraph{Marginal likelihood tuning.}
The prior precision $\tau$ is selected by maximising the log evidence:
\begin{equation}
  \log p(\cD)
    \approx \log p(\cD|\btheta_{\MAP})
    + \log p(\btheta_{\MAP})
    + \tfrac{P}{2}\log 2\pi
    - \tfrac{1}{2}\log\det\mathbf{H},
  \label{eq:marglik}
\end{equation}
which is differentiable with respect to $\tau$ and requires no validation
data~\citep{daxberger2021laplace}. Grid search over
a discrete set of values on validation NLL is simpler but requires a
held out split.

\section{Evaluation Metrics}
\label{sec:metrics}

\paragraph{Brier Score.}
$\mathrm{BS}=\frac{1}{N}\sum_n\sum_c(\hat{p}_{nc}-\mathbf{1}[y_n=c])^2$;
bounded and less sensitive to extreme probabilities than NLL.

\paragraph{Mean Confidence / Mean Entropy / $1 - \text{Confidence}$.}
Mean max class probability, mean Shannon entropy $H$, and
$1-\max_c\hat{p}_c\in[0,0.5]$.

\paragraph{Uncertainty calibration U MAE, U MSE, U Corr, Max Gap.}
Compares expected uncertainty against empirical error rate
per bin. Pearson correlation and max bin gap assess worst case alignment.

\paragraph{Epistemic uncertainty OOD maps.}
Mutual information $I(y;\btheta|\bm{x},\cD)$ estimated via
equation~\eqref{eq:mi} with $S=50$ MC samples.

\section{Shared Infrastructure}
\label{sec:shared}

Both notebooks share a \texttt{shared} utility package. The MAP baseline uses Adam with early stopping.

\section{Two Moons Classification}
\label{sec:twomoons}

\subsection{Dataset and OOD Probe Design}

2000 2D points from two interleaved half crescents.
To probe out of distribution behaviour, 200 probe points are
sampled in four corner clusters deliberately outside the training manifold.

\subsection{\texttt{bayesian-torch}: Variational BNN}
\label{sec:bnn}

\begin{lstlisting}[caption={BNN conversion and ELBO AvUC training.}]
dnn_to_bnn(net, {"prior_mu": 0.0, "prior_sigma": 1.0,
    "posterior_mu_init": 0.0, "posterior_rho_init": -3.0,
    "type": "Reparameterization",
    "moped_enable": True, "moped_delta": 0.5})

loss = CE(mean_logits, labels) + get_kl_loss(net) / N + avu
\end{lstlisting}

Table~\ref{tab:bnn_main_subset} reports the relevant calibration statistics for all \texttt{bayesian-torch} configurations.

\begin{table}[H]
  \centering
  \caption{Relevant calibration metrics for bayesian torch variants on Two Moons.}
  \label{tab:bnn_main_subset}
  \small\setlength{\tabcolsep}{4pt}
  \begin{tabular}{llccccccc}
    \toprule
    Type & Init & Loss & BS & $1 - c$ & U MAE & U MSE & U Corr & MaxGap \\
    \midrule
    Flipout & scr & ELBO & 0.058 & 0.094 & 0.018 & 0.458 & 0.988 & 0.069 \\
    Flipout & scr & AvUC & 0.059 & 0.072 & 0.014 & 0.840 & 0.953 & 0.134 \\
    Flipout & MOP & ELBO & 0.059 & 0.089 & 0.019 & 0.882 & 0.935 & 0.113 \\
    Flipout & MOP & AvUC & 0.059 & 0.072 & 0.010 & 0.294 & 0.990 & 0.078 \\
    BBB & scr & ELBO & 0.058 & 0.088 & 0.019 & 0.582 & 0.983 & 0.051 \\
    BBB & scr & AvUC & 0.095 & 0.153 & 0.028 & 0.918 & 0.965 & 0.102 \\
    BBB & MOP & ELBO & 0.058 & 0.087 & 0.015 & 0.556 & 0.967 & 0.130 \\
    BBB & MOP & AvUC & 0.057 & 0.067 & 0.014 & 0.900 & 0.941 & 0.123 \\
    \bottomrule
  \end{tabular}
\end{table}

\begin{table}[H]
  \centering
  \caption{OOD epistemic uncertainty decomposition for all bayesian torch variants.}
  \label{tab:ood}
  \small\setlength{\tabcolsep}{4pt}
  \begin{tabular}{lccccccc}
    \toprule
    Variant & ID $1 - c$ & OOD $1 - c$ & ID Alea & OOD Alea & ID Epist & OOD Epist & $\Delta$Epist \\
    \midrule
    Flip scr ELBO & 0.094 & 0.005 & 0.238 & 0.028 & 0.003 & 0.001 & $-0.003$ \\
    Flip scr AvUC & 0.072 & 0.003 & 0.177 & 0.012 & 0.008 & 0.001 & $-0.007$ \\
    Flip MOP ELBO & 0.089 & 0.010 & 0.213 & 0.034 & 0.009 & 0.003 & $-0.006$ \\
    Flip MOP AvUC & 0.073 & 0.004 & 0.150 & 0.018 & 0.036 & 0.004 & $-0.032$ \\
    BBB scr ELBO  & 0.089 & 0.003 & 0.227 & 0.018 & 0.002 & 0.000 & $-0.002$ \\
    BBB scr AvUC  & 0.153 & 0.032 & 0.359 & 0.123 & 0.003 & 0.002 & $-0.001$ \\
    BBB MOP ELBO  & 0.086 & 0.011 & 0.202 & 0.039 & 0.014 & 0.004 & $-0.009$ \\
    BBB MOP AvUC  & 0.067 & 0.007 & 0.140 & 0.027 & 0.035 & 0.009 & $-0.026$ \\
    \bottomrule
  \end{tabular}
\end{table}

Please refer to the Appendix for histograms, boundary visualizations, and full tables.

\subsection{\texttt{laplace-torch}: Laplace Approximation}
\label{sec:la}

\begin{lstlisting}[caption={Laplace post hoc workflow.}]
la = Laplace(model, 'classification',
             subset_of_weights='last_layer', hessian_structure='kron')
la.fit(train_loader)
la.optimize_prior_precision(method='gridsearch', val_loader=val_loader,
                            pred_type='glm', link_approx='probit')
probs = la(X_test, pred_type='glm', link_approx='probit')
\end{lstlisting}

Table~\ref{tab:la_main_subset} reports the relevant subset of metrics for the requested Laplace configurations including Diagonal all weights, Full Last Layer, and Kronecker Last Layer under both Marginal Likelihood and Gridsearch tuning.

\begin{table}[H]
  \centering
  \caption{Relevant calibration metrics for laplace torch on Two Moons.}
  \label{tab:la_main_subset}
  \small\setlength{\tabcolsep}{4pt}
  \begin{tabular}{llccccccc}
    \toprule
    Variant & Scope & PP & BS & $1 - c$ & U MAE & U MSE & U Corr & MaxGap \\
    \midrule
    Diag all GS & all & $10^4$ & 0.056 & 0.079 & 0.010 & 0.806 & 0.959 & 0.197 \\
    Full LL GS  & LL  & $10^4$ & 0.056 & 0.079 & 0.010 & 0.806 & 0.959 & 0.197 \\
    Kron LL GS  & LL  & $10^4$ & 0.056 & 0.079 & 0.010 & 0.806 & 0.959 & 0.197 \\
    Diag all ML & all & 8.73  & 0.057 & 0.088 & 0.016 & 0.817 & 0.955 & 0.196 \\
    Full LL ML  & LL  & 1.79  & 0.062 & 0.145 & 0.070 & 6.49  & 0.935 & 0.136 \\
    Kron LL ML  & LL  & 2.04  & 0.061 & 0.138 & 0.064 & 5.50  & 0.938 & 0.141 \\
    \bottomrule
  \end{tabular}
\end{table}

Please refer to the Appendix for epistemic maps, full metrics, and visualizations.

\section{Discussion}
\label{sec:discussion}

MOPED is a geometric anchor. The BBB scratch AvUC failure demonstrates that AvUC calibration alone is not sufficient.
Under marglik, full matrix GGN curvature is poorly conditioned and selects pathological prior precisions.
All weights LA marglik variants retain some OOD epistemic sensitivity, while last layer LA and all SVI variants collapse to near zero outside the boundary.

\appendix
\section{Appendix: Full Metric Tables and Visualizations}

\begin{sidewaystable}
  \centering
  \caption{All calibration metrics for bayesian torch variants on Two Moons.}
  \label{tab:bnn_main_full}
  \scriptsize\setlength{\tabcolsep}{2.5pt}
  \begin{tabular}{
    L{1.5cm} C{0.6cm} C{0.65cm} C{0.65cm} C{0.7cm} C{0.7cm} C{0.7cm}
    C{0.7cm} C{0.7cm} C{0.65cm} C{0.7cm} C{0.95cm} C{0.7cm}
    C{0.75cm} C{0.7cm} C{0.75cm} C{0.5cm} C{0.9cm} C{0.85cm}
  }
    \toprule
    Type & Init & Loss & Acc & NLL & BS & ECE & Conf & Ent & $1{-}c$
    & U MAE & U MSE & U Corr & MaxGap & ClECE & MActErr & Ep & Fit & Inf \\
    \midrule
    Flipout & scr & ELBO & 0.921 & 0.198 & 0.058 & 0.018 & 0.906 & 0.241 & 0.094 & 0.018 & 0.458 & 0.988 & 0.069 & 0.018 & 0.080 & 39 & 19.1 & 0.064 \\
    Flipout & scr & AvUC & 0.918 & 0.198 & 0.059 & 0.014 & 0.928 & 0.184 & 0.072 & 0.014 & 0.840 & 0.953 & 0.134 & 0.014 & 0.082 & 38 & 80.4 & 0.065 \\
    Flipout & MOP & ELBO & 0.918 & 0.199 & 0.059 & 0.019 & 0.911 & 0.222 & 0.089 & 0.019 & 0.882 & 0.935 & 0.113 & 0.019 & 0.083 & 39 & 20.0 & 0.063 \\
    Flipout & MOP & AvUC & 0.920 & 0.201 & 0.059 & 0.010 & 0.928 & 0.184 & 0.072 & 0.010 & 0.294 & 0.990 & 0.078 & 0.010 & 0.081 & 40 & 90.1 & 0.064 \\
    BBB & scr & ELBO & 0.923 & 0.197 & 0.058 & 0.019 & 0.911 & 0.229 & 0.088 & 0.019 & 0.582 & 0.983 & 0.051 & 0.019 & 0.078 & 50 & 23.2 & 0.051 \\
    BBB & scr & AvUC & 0.866 & 0.311 & 0.095 & 0.028 & 0.847 & 0.362 & 0.153 & 0.028 & 0.918 & 0.965 & 0.102 & 0.028 & 0.134 & 32 & 73.3 & 0.050 \\
    BBB & MOP & ELBO & 0.920 & 0.195 & 0.058 & 0.015 & 0.913 & 0.217 & 0.087 & 0.015 & 0.556 & 0.967 & 0.130 & 0.015 & 0.081 & 36 & 15.9 & 0.055 \\
    BBB & MOP & AvUC & 0.923 & 0.192 & 0.057 & 0.014 & 0.933 & 0.173 & 0.067 & 0.014 & 0.900 & 0.941 & 0.123 & 0.014 & 0.078 & 33 & 68.7 & 0.051 \\
    \bottomrule
  \end{tabular}
\end{sidewaystable}

\begin{sidewaystable}
  \centering
  \caption{All calibration metrics for laplace torch variants on Two Moons.}
  \label{tab:la_main_full}
  \scriptsize\setlength{\tabcolsep}{2.5pt}
  \begin{tabular}{
    L{1.9cm} C{0.6cm} C{0.55cm} C{0.7cm} C{0.7cm} C{0.7cm} C{0.7cm}
    C{0.7cm} C{0.7cm} C{0.65cm} C{0.7cm} C{0.95cm} C{0.7cm}
    C{0.75cm} C{0.7cm} C{0.95cm} C{0.75cm} C{0.95cm}
  }
    \toprule
    Variant & Hess & Sc & PP & NLL & BS & ECE & Conf & Ent & $1{-}c$
    & U MAE & U MSE & U Corr & MaxGap & ClECE & Fit & Tune & Inf \\
    \midrule
    Diag all GS & diag & all & $10^4$ & 0.190 & 0.056 & 0.010 & 0.921 & 0.194 & 0.079 & 0.010 & 0.806 & 0.959 & 0.197 & 0.010 & 0.23  & 4.01 & 0.021 \\
    Full all GS & full & all & $10^4$ & 0.190 & 0.056 & 0.010 & 0.921 & 0.194 & 0.079 & 0.010 & 0.806 & 0.959 & 0.197 & 0.010 & 1263  & 11.4 & 0.034 \\
    Kron all GS & kron & all & $10^4$ & 0.190 & 0.056 & 0.010 & 0.921 & 0.194 & 0.079 & 0.010 & 0.806 & 0.959 & 0.197 & 0.010 & 0.54  & 5.29 & 0.022 \\
    Full LL GS  & full & LL  & $10^4$ & 0.190 & 0.056 & 0.010 & 0.921 & 0.194 & 0.079 & 0.010 & 0.806 & 0.959 & 0.197 & 0.010 & 7.12  & 1.52 & 0.001 \\
    Kron LL GS  & kron & LL  & $10^4$ & 0.190 & 0.056 & 0.010 & 0.921 & 0.194 & 0.079 & 0.010 & 0.806 & 0.959 & 0.197 & 0.010 & 0.43  & 1.85 & 0.001 \\
    Diag all ML & diag & all & 8.73  & 0.192 & 0.057 & 0.016 & 0.912 & 0.219 & 0.088 & 0.016 & 0.817 & 0.955 & 0.196 & 0.016 & 0.20  & 0.06 & 0.021 \\
    Full all ML & full & all & 0.52  & 0.307 & 0.082 & 0.152 & 0.771 & 0.500 & 0.229 & 0.152 & 26.0  & 0.913 & 0.218 & 0.152 & 1122  & 1.88 & 0.034 \\
    Kron all ML & kron & all & 1.42  & 0.238 & 0.065 & 0.086 & 0.841 & 0.377 & 0.159 & 0.086 & 9.07  & 0.923 & 0.152 & 0.086 & 0.53  & 0.10 & 0.022 \\
    Full LL ML  & full & LL  & 1.79  & 0.226 & 0.062 & 0.070 & 0.855 & 0.348 & 0.145 & 0.070 & 6.49  & 0.935 & 0.136 & 0.070 & 7.13  & 0.09 & 0.001 \\
    Kron LL ML  & kron & LL  & 2.04  & 0.221 & 0.061 & 0.064 & 0.862 & 0.334 & 0.138 & 0.064 & 5.50  & 0.938 & 0.141 & 0.064 & 0.44  & 0.08 & 0.001 \\
    \bottomrule
  \end{tabular}
\end{sidewaystable}

\begin{figure}[H]
  \centering
  \includegraphics[width=0.60\linewidth]{ood_points.png}
  \caption{Training data and OOD probe points.}
  \label{fig:ood_points}
\end{figure}

\begin{figure}[H]
  \centering
  \begin{subfigure}{0.48\linewidth}
    \includegraphics[width=\linewidth]{bbb_moped_avuc_boundary.png}
    \caption{BBB MOPED AvUC boundary}
  \end{subfigure}\hfill
  \begin{subfigure}{0.48\linewidth}
    \includegraphics[width=\linewidth]{bbb_scratch_avuc_boundary.png}
    \caption{BBB scratch AvUC boundary}
  \end{subfigure}
  \\[6pt]
  \begin{subfigure}{0.48\linewidth}
    \includegraphics[width=\linewidth]{bbb_moped_avuc_uncertainty.png}
    \caption{BBB MOPED AvUC uncertainty map}
  \end{subfigure}\hfill
  \begin{subfigure}{0.48\linewidth}
    \includegraphics[width=\linewidth]{bbb_scratch_avuc_uncertainty.png}
    \caption{BBB scratch AvUC uncertainty}
  \end{subfigure}
  \caption{MOPED vs scratch initialisation with AvUC loss.}
  \label{fig:bnn_boundaries}
\end{figure}

\begin{figure}[H]
  \centering
  \begin{subfigure}{0.48\linewidth}
    \includegraphics[width=\linewidth]{bbb_moped_avuc_histogram.png}
    \caption{BBB MOPED AvUC histogram}
  \end{subfigure}\hfill
  \begin{subfigure}{0.48\linewidth}
    \includegraphics[width=\linewidth]{bbb_moped_avuc_reliability.png}
    \caption{BBB MOPED AvUC reliability}
  \end{subfigure}
  \\[6pt]
  \begin{subfigure}{0.48\linewidth}
    \includegraphics[width=\linewidth]{bbb_scratch_avuc_histogram.png}
    \caption{BBB scratch AvUC histogram}
  \end{subfigure}\hfill
  \begin{subfigure}{0.48\linewidth}
    \includegraphics[width=\linewidth]{bbb_moped_elbo_reliability.png}
    \caption{BBB MOPED ELBO reliability}
  \end{subfigure}
  \caption{Probability histograms and reliability diagrams.}
  \label{fig:bnn_histrel}
\end{figure}

\begin{figure}[H]
  \centering
  \begin{subfigure}{0.48\linewidth}
    \includegraphics[width=\linewidth]{bbb_moped_avuc_ood_epistemic.png}
    \caption{BBB MOPED AvUC}
  \end{subfigure}\hfill
  \begin{subfigure}{0.48\linewidth}
    \includegraphics[width=\linewidth]{bbb_scratch_elbo_ood_epistemic.png}
    \caption{BBB scratch ELBO}
  \end{subfigure}
  \caption{OOD epistemic uncertainty.}
  \label{fig:bnn_ood}
\end{figure}

\begin{figure}[H]
  \centering
  \begin{subfigure}{0.48\linewidth}
    \includegraphics[width=\linewidth]{marglik_full_all_weights_boundary.png}
    \caption{Full all marglik boundary}
  \end{subfigure}\hfill
  \begin{subfigure}{0.48\linewidth}
    \includegraphics[width=\linewidth]{marglik_full_all_weights_reliability.png}
    \caption{Full all marglik reliability}
  \end{subfigure}
  \caption{Marginal likelihood failure for Full all LA.}
  \label{fig:la_failure}
\end{figure}

\begin{figure}[H]
  \centering
  \begin{subfigure}{0.48\linewidth}
    \includegraphics[width=\linewidth]{marglik_kron_last_layer_uncertainty.png}
    \caption{Kron LL marglik}
  \end{subfigure}\hfill
  \begin{subfigure}{0.48\linewidth}
    \includegraphics[width=\linewidth]{marglik_kron_all_weights_ood_epistemic.png}
    \caption{Kron all marglik OOD epistemic}
  \end{subfigure}
  \caption{Uncertainty patterns under marginal likelihood tuning.}
  \label{fig:la_stripe}
\end{figure}

\bibliographystyle{plainnat}
\bibliography{references}

\end{document}